% SPDX-License-Identifier: CC-BY-4.0
% Copyright (C) 2026 Fabian Januszewski, Abhijit (the Langlandscape team)
% GENERATED by tools/extract_statements.py from mathematics/chapters/ -- do not edit by hand.

\begin{quoted}{Hypothesis}{20.1}{165}{hyp:s06-H}{standing hypotheses; seed \S9.1, corrected}{chapters/s06/s06-m2.tex}
\leavevmode
\begin{itemize}
\item[(H0)] $p\ge5$, $q\ne p$, $(M,pq)=1$, $N=qM$, $p\nmid\varphi(N)$. In particular $q\not\equiv1\bmod p$.
\item[(H0$^-$)] (\emph{new}; needed for the fixed-weight input Cor.~21.2 and for Prop.~24.1;
  Thm~21.5 needs (H0$^-$) \emph{or} (Lev); the inclusion of Prop.~22.1 uses neither)
  $q\not\equiv-1\bmod p$, or $\bar\rho_\m$ is ramified at $q$. Together with (H0): $p\nmid q^2-1$, or $\bar\rho_\m$ is ramified at $q$.
\item[(Lev)] (\emph{new}; alternative to (H0$^-$) in Thm~21.5) At the point $P$ of (P), of weight $k=k_P$, with
  $f=F_P$ (ordinary): either (F0) $M=1$; or (F1) $M>1$ is squarefree, $\bar\rho$ is unramified at every $\ell\mid M$ (so $N(\bar\rho)=1$),
  and $\ell\not\equiv-1\bmod p$ for every $\ell\mid M$; or (F1$'$) as (F1), but with some $\ell\mid M$, $\ell\equiv-1\bmod p$. These are the
  cases (F0), (F1), (F1$'$) of Prop.~23.2, in which its hypothesis (Lev$'_M$) (with (L3)) holds automatically. Status:
  Prop.~23.2 is \stpmc{} in all three cases, for (F1$'$) in our sign-fixed reading R4 of \cite[Thm~3.14]{KO} (ours, not
  \cite{KO}'s text). Not included: $N(\bar\rho)>1$ and non-squarefree $M_0$, which are \stopen{} (Remark~23.8); (L3) is the condition under
  which the method may extend.

\item[(H1)] $\T$ is Hida's full ordinary cuspidal Hecke algebra of tame level $\Gamma_0(N)$ over $\Lam$, in a fixed weight component
  $k_0\bmod(p-1)$. It is generated by all $T_\ell$ ($\ell\nmid Np$), $U_\ell$ ($\ell\mid Np$, in particular $U_q$ and $U_p$), and it
  is localised at a non-Eisenstein maximal ideal $\m$. $F$ is a primitive ordinary $q$-new branch, with $\lambda_F:\T\to I_F$, and
  $I_F$ is a normal local domain finite over $\Lam$. Then $I_F$ is $\Lam$-free: it is normal of dimension $2$, hence
  Cohen--Macaulay (Serre's criterion), hence free over the regular local ring $\Lam$ (Auslander--Buchsbaum). So the seed's
  ``hence $\Lam$-free'' is correct; it is used only through flatness of $R=\T\otimes_\Lam I_F$ over $I_F$.
\item[(H2)] $\bar\rho=\bar\rho_\m$ is absolutely irreducible on $G_{\Q(\sqrt{p^*})}$ and $p$-distinguished.
\item[(Im)] $p\mid\#\bar\rho(G_\Q)$. \emph{None of the cited theorems below uses (Im).} \cite[Thm~1.1]{KO} needs only (H2) (its
  condition (2) $=$ Assumption 2.1(TW)), and \cite[Rem.~2.10]{KO} says that the \cite{DFG} and \cite{Dim09} image conditions are interchangeable
  over $\Q$. We keep (Im) as in the seed, flagged as probably superfluous (Risk R4).
\item[(P)] Either (P1) $k_0\not\equiv2\bmod(p-1)$, and then $F$ has an arithmetic point $P$ of weight $k_P\in[3,p-1]$ with trivial
  character (a \emph{Kim--Ota point}). Or (P2) $k_0\equiv2$, $F$ has a weight-two point $P$ which is the ordinary
  $p$-stabilisation of a newform $f$ of level $N$, \emph{and} (\emph{new}) the weight-two ordinary Hecke algebra
  $\T^{\ord}_2(\Gamma_0(Np))_\m$ has no $p$-new eigenforms. The extra condition makes $\T/P\T\cong\T_2(\Gamma_0(N))_\m$,
  which is the algebra of \cite{KO}. Without it the control step (3) of the proof compares with an algebra that \cite{KO} does not treat.
\end{itemize}
\sthyp
\end{quoted}

\begin{quoted}{Theorem}{19.8}{163}{thm:s06-E}{the combinatorial half; seed Thm~E, hypotheses explicit}{chapters/s06/s06-m1.tex}
Let $S,\fa,\fb$ be as in Lemma~19.5. Put $S'=S/\fa$, and let $\lambda:S\to A$ factor through $\lambda':S'\to A$. Assume
\begin{itemize}
\item[(Gor$_S$)] $S$ is $A$-Gorenstein.
\end{itemize}
Let $\theta$ generate $S[\lambda]$ (Lemma~19.3), and let $x\in S'$ be its image. Then:
\begin{enumerate}[label=(\roman*)]
\item $\theta\in\fb$, $x\in S'[\lambda']$, $\eta(S,\lambda)=(\lambda'(x))$, and $\cc_q:=\lambda(\fb)=\{\psi(x):\psi\in S'^\vee\}$ (the
  \emph{order ideal} of $x$).
\item If $S'[\lambda']$ is free of rank one with a generator $\theta'$ that is unimodular in $S'$ (i.e.\ $S'[\lambda']$ is an
  $A$-direct summand of $S'$), then $\eta(S,\lambda)=\cc_q\cdot\eta(S',\lambda')$.
\item This holds in particular under (Gor$_{S'}$): $S'$ is finite free and $A$-Gorenstein. It also holds whenever $A$ is a DVR
  (\emph{no} hypothesis on $S'$). For $A$ normal, the identity therefore holds after localisation at every height-one prime.
\item In general $\cc=\mathrm{im}(\fb\to S')\cong\fb\cong S'^\vee$. So, if $S'$ is $A$-free (e.g.\ $S/\fb$ free, Lemma~19.5(c)), $\cc$ is a principal ideal of $S'$ if and only if
  (Gor$_{S'}$) holds.
\end{enumerate}
With $S=R$, $A=I_F$: under (Gor$_R$) and (Gor$_{R'}$), $\eta_F=\cc_q(F)\cdot\eta^B_F$, i.e.\ $(h_F)=\cc_q(F)$. \stproved{}
(Lean, namespace \pathtt{Langlandscape.CommAlg.ThmE}: (i) \lean{Langlandscape.CommAlg.ThmE.theta\_mem\_b}, \pathtt{mk\_theta\_mem\_ann}, \pathtt{eta\_eq\_span\_lambda'}, \lean{Langlandscape.CommAlg.ThmE.c\_eq\_orderIdeal};
(ii) \lean{Langlandscape.CommAlg.ThmE.eta\_eq\_c\_mul\_of\_unimodular};
(iii) Gorenstein half \lean{Langlandscape.CommAlg.ThmE.unimodular\_of\_isGorenstein}, \lean{Langlandscape.CommAlg.ThmE.eta\_eq\_c\_mul\_of\_isGorenstein}, DVR half \lean{Langlandscape.CommAlg.ThmE.eta\_eq\_c\_mul\_of\_dvr} (\pathtt{CommAlg/DVR.lean}, with \pathtt{exists\_unimodular\_of\_dvr});
(iv) \lean{Langlandscape.CommAlg.ThmE.c\_principal\_iff\_isGorenstein}.)
\stformalized{} for (i), (ii), (iii) (both halves) and (iv) for $S'$ $A$-free, under the hypotheses of Lemma~19.5 as the structure \texttt{NewOldData} (same qualifiers as there). \emph{Not formalized:} the localisation clause of (iii) for $A$ normal (height-one primes), the isomorphism $\mathrm{im}(\fb\to S')\cong\fb\cong S'^\vee$ as a standalone statement for general $S'$, and the final specialisation $S=R$, $A=I_F$ (the abstract statement is formalized, its instantiation to the Hecke rings is not).
\end{quoted}

\begin{quoted}{Proposition}{22.1}{174}{prop:s06-galois}{Galois comparison; seed Prop.~galois}{chapters/s06/s06-m4.tex}
Assume (H0), (H1), (H2). Then $\tq_q(F)\subseteq\cc_q(F)$. \stpmc{}
(Lean: Given \pathtt{Hyp.ResidualBurnside}, \pathtt{Hyp.FreeRankTwo} is equivalent to the inclusion $\tq_q\subseteq\cc_q$; both directions are proved in Lean (\pathtt{Defect.galois}, \pathtt{Defect.freeRankTwo\_of\_le}, packaged as \pathtt{Defect.freeRankTwo\_iff\_le}). The Galois half is therefore not formalized independently: its arithmetic content (that the old-form representation $\rho_M$ exists at level $M$, steps 1--3 of Prop.~22.1) is the cited hypothesis. \pathtt{Hyp.ResidualBurnside} is discharged only given absolute irreducibility, (H2).)
\par\smallskip\noindent{\scriptsize\itshape [Status paragraph following the statement in the chapter:]}\par
\noindent\emph{Cited inputs.}
(1) $H$ is free of rank two over $\T$, and $H(M)$ is free of rank two over $\T(M)_\m$ (tame level $M$):
\cite[\S7.2, p.~68]{Hs} \lit{hsieh\_2017\_hida-families-and-p-adic-triple-product-l-functions\_arxiv}{\S7.2, proof of Lemma 7.4, p.~68}
\stcited, citing \cite[Cor.~1, 2, p.~482]{Wi95}
\lit{wiles\_1995\_modular-elliptic-curves-and-fermats-last-theorem}{Thm 2.1, Cor.~1, Cor.~2, p.~482}, under (CR) (which (H2)
implies, Remark~20.2). The Galois action on $H$ (resp.\ $H(M)$) is unramified outside $Np$ (resp.\ $Mp$), and
$\tr(\Frob_\ell\mid H)=T_\ell$ in $\T$ for $\ell\nmid Np$ (Eichler--Shimura; for each classical specialisation this is
\lit{darmon-diamond-taylor\_1997\_fermats-last-theorem\_preprint}{Thm 3.1(a), p.~85}) \stcited.
(2) $H(M)$ is unramified at $q$, since $q\nmid Mp$ ($X_1(Mp^r)$ has good reduction at $q$; for each classical specialisation this is
again \lit{darmon-diamond-taylor\_1997\_fermats-last-theorem\_preprint}{Thm 3.1(a), p.~85}) \stknown.
(3) Chebotarev: for a Galois extension of $\Q$ unramified outside a finite set $S$, the Frobenius elements at primes $\notin S$ are
dense \lit{darmon-diamond-taylor\_1997\_fermats-last-theorem\_preprint}{Thm 2.3, p.~52} \stknown.
\end{quoted}

\begin{quoted}{Theorem}{21.5}{172}{thm:s06-defect}{$\Lam$-adic defect formula; seed Thm~defect, corrected}{chapters/s06/s06-m3.tex}
Assume (H0), (H1), (H2), (P) (with (P2) in its corrected form), and (H0$^-$) or (Lev) (Hyp.~20.1). Then
\[(h_F)=\cc_q(F)=\tq_q(F)\qquad\text{as principal ideals of }I_F.\]
\stpmc{} under (H0$^-$). Under (Lev) (this is where $q\equiv-1\bmod p$ with $\bar\rho$ unramified at $q$ is covered): \stpmc{} for $M=1$
and for $M>1$ in the cases (F1), (F1$'$) of Prop.~23.2 (for (F1$'$) in our sign-fixed reading of \cite[Thm~3.14]{KO}, not
\cite{KO}'s text). For $q\equiv-1\bmod p$ with $\bar\rho$ unramified at $q$ and outside (Lev) ($N(\bar\rho)>1$, or $\ell^2\mid M$ at a prime $\ell$ where $\bar\rho$ is unramified) the theorem is \stopen{} (Remark~25.1(i), Remark~23.8).

 Cited inputs: Hida freeness
and control \cite[Thm~3.1, Cor.~3.2]{Hi86b} \lit{hida\_1986\_iwasawa-modules-attached-to-congruences-of-cusp-forms}{Thm 3.1, Cor.~3.2
(3.3b), p.~250}; (Gor$_R$) and (Gor$_{R'}$), with freeness of $H$ (\cite{Hs} as above, resting on \cite{Wi95}, \cite{MW86}, \cite[Cor.~8.11]{Hel}
\lit{helm\_2004\_maps-between-modular-jacobians-and-jacobians-of-shimura-curves\_arxiv}{Cor.~8.11, p.~31}); the $\Lam$-adic JL
(\pathtt{if:s01-targets}); and Thm~21.1 at the point $P$ of (P) under (H0$^-$), or Prop.~23.2 at $P$ under
(Lev) (resting on \cite[Thm~2.9, Lemma~2.4]{KO} for $M=1$, \cite[Thm~3.14]{KO} for (F1), (F1$'$)). (Im) is not used.
\end{quoted}

\begin{quoted}{Corollary}{21.6}{173}{cor:s06-allweights}{all weights of the branch; seed Cor.~allweights}{chapters/s06/s06-m3.tex}
Under the hypotheses of Thm~21.5, let $P'$ be any arithmetic point of weight $k'\ge3$ with trivial character
(including $k'\ge p$), with $q$-new $f'=F_{P'}$ ($F_{P'}$ is the $p$-stabilisation of a newform of level $N$, $F$ being primitive of tame level $N$). Then $\cc_q(f')=(\varpi_{P'}^{t_{f'}(q)})$ (in $\widetilde\cO_{P'}$) and
$\ord_{\varpi_{P'}}\eta_{f'}(N)-\ord_{\varpi_{P'}}\eta_{f'}(M,q)=t_{f'}(q)$, for the ring-theoretic congruence ideals of $\T_{k'}(N)_\m$ and of its
$q$-new quotient. \stpmc{} (with the status of the case of Thm~21.5 that is used; under (Lev), case (F1$'$), in our reading of \cite[Thm~3.14]{KO}).
\end{quoted}

\begin{quoted}{Proposition}{23.2}{176}{prop:s06-levM}{the inclusion without (H0$^-$); the seed's step (4)}{chapters/s06/s06-m5.tex}
Let $f$ be as above, with (H0) and (Lev$'_M$) (Hyp.~23.1), and assume that $f$ is $p$-ordinary, i.e.\ $a_p(f)\in\cO^\times$
(automatic for $f=F_P$ in Thm~21.5). Then $\cc_q(f)\subseteq(\varpi^{t_f(q)})$ in each of the following cases, in which
(Lev$'_M$) holds automatically, $N(\bar\rho)=1$, $M_0=M$ and $\Sigma^+=P_{\bar\rho}=\emptyset$:
\begin{itemize}
\item[(F0)] $M=1$. \stpmc{} (Remark~23.10).
\item[(F1)] $M>1$ squarefree, $\bar\rho$ unramified at every $\ell\mid M$, and $\ell\not\equiv-1\bmod p$ for every $\ell\mid M$.
  \stpmc
\item[(F1$'$)] as (F1), but some $\ell\mid M$ has $\ell\equiv-1\bmod p$. \stpmc{} in the sign-fixed reading of \cite[Thm~3.14]{KO}
  (Remark~23.4), which is our reading, not \cite{KO}'s text (F2 part B and audit part 2: confirmed as a reading).

\end{itemize}
No condition on $q\bmod p$ beyond (H0) is used; in particular $q\equiv-1\bmod p$ is allowed. For $N(\bar\rho)>1$ (under (Lev$'_M$), with
or without (L3)) the inclusion is \stopen: the proof below does not apply when $\bar\rho$ is not twist-minimal (counterexample B2 of F2,
Remark~23.9), and the inclusion itself is neither proved nor disproved there.
\end{quoted}

\begin{quoted}{Proposition}{24.1}{184}{prop:s06-betti}{Betti incarnation of $\cc_q$}{chapters/s06/s06-m7.tex}
Let $l=p>k\ge4$ (the range of \pathtt{if:s03-amalgam}), $q\not\equiv-1\bmod p$, and let $\m$ be non-Eisenstein for $\T_k(M)$ acting on
$H^1(\Gamma_0(M),V_k)$, with $H^1(\Gamma_0(M),V_k)_\m$ free of rank two over $\T_k(M)_\m$. Then:
\begin{enumerate}[label=(\alph*)]
\item $\coker(\Psi^*\Psi)_\m\cong\big(\T_k(M)_\m/(\delta)\big)^{\oplus2}$, where $\delta:=\det(\Psi^*\Psi)=(q+1)^2-q^{2-k}T_q^2$.
  Since $p\nmid2(q+1)$, $(\delta)=(T_q-\epsilon(q+1)q^{(k-2)/2})$ for the unique sign $\epsilon$ with this element in $\m$.
\item $\T_k(N)_{\m}/(\fa+\fb)\cong\T_k(M)_\m/(T_q-\epsilon(q+1)q^{(k-2)/2})$ (Ribet's level-raising module, with Ihara's lemma).
\item If $\rk_\cO\T'_{k,\m}=1$, then $\ord_\varpi\cc_q(f)=\tfrac12\,\mathrm{length}_\cO\coker(\Psi^*\Psi)_\m$. In general only
  $\ord_\varpi\cc_q(f)\le\tfrac12\mathrm{length}$ holds.
\end{enumerate}
(a) \stproved{} (linear algebra, Lemma~24.2). (b) \stpmc{} (Ihara's lemma in weight $k$ and the cohomological
congruence module: \cite[Lemma~3.2 and \S2]{Dia91}, \cite[\S3]{Dia89}). (c) \stpmc{} (from (a), (b)).
Here: in (a), if neither sign gives an element of $\m$, then $\delta\in\T_k(M)_\m^\times$ and both sides are $0$ for either sign;
in (b), $\T_k(N)_\m$ is the localisation of the weight-$k$ Hecke algebra of $\Gamma_0(N)$ (with $U_q$; $N=qM$) at the maximal ideal $\m_N$
generated by $\m$ (on the $T_\ell$, $\ell\nmid Np$, and $U_\ell$, $\ell\mid M$) and $U_q-\epsilon c$, and $\fa,\fb$ are its $q$-new and
$q$-old kernels; in (c), $f$ is a $q$-new eigenform in $\m_N$, $\T'_{k,\m}$ is the $q$-new quotient of $\T_k(N)_\m$, $\cc_q(f)$ is as in
Def.~20.5, and (H0) is assumed. \provisional{S06-M7: the degenerate case ``no sign'' in (a) and the maximal ideal
$\m_N$ of (b) are made explicit, and (H0) ($p\nmid\varphi(N)$, $q\not\equiv1$) is used in (b); statement otherwise unchanged.}
\end{quoted}

\begin{quoted}{Conjecture}{25.2}{187}{conj:s06-several}{several ramified primes; seed \S10, last paragraph}{chapters/s06/s06-m8.tex}
Let $D=q_1\cdots q_r$ ($r$ odd), with each $q_i$ satisfying (H0), (H0$^-$) and $F$ $D$-new. Let
$\T=\T^{(0)}\twoheadrightarrow\T^{(1)}\twoheadrightarrow\cdots\twoheadrightarrow\T^{(r)}=\T^{D\text{-new}}\cong\T^{\ord}_{B_D}$ be the
$q_1\cdots q_i$-new quotients, and let $\cc^{(i)}_{q_i}(F)$ be the old/new congruence ideal of $\T^{(i-1)}\twoheadrightarrow\T^{(i)}$ at
$q_i$. Then $(h_F)=\prod_i\cc^{(i)}_{q_i}(F)=\prod_i\tq_{q_i}(F)$. \stconj{} If every $\T^{(i)}_\m$ is Gorenstein, the first
equality follows by iterating Thm~19.8 \stsketch. Whether the intermediate quotients are Gorenstein is \stopen{} (Hsieh
covers $i=0$ and $i=r$). The second equality would follow from Prop.~22.1 at each step and \cite[Thm~1.1]{KO} with
$N^-=D$, by the squeeze of step (5) \stsketch. Note that $q_i$ is $q_1\cdots q_{i-1}$-new-ramified in the intermediate algebras,
which needs the relative version of Lemma~20.4 at each step.
\end{quoted}
